\documentclass[11pt]{article}

\usepackage[margin=1in]{geometry}
\usepackage{amsmath,amssymb,amsthm}
\usepackage{graphicx}
\usepackage[hidelinks]{hyperref}
\usepackage{enumitem}
\usepackage{array}

\setlength{\parindent}{0pt}
\setlength{\parskip}{0.65em}
\setlist[itemize]{leftmargin=1.5em}

\newtheorem{theorem}{Theorem}
\newtheorem{lemma}{Lemma}

\newcommand{\R}{\mathbb{R}}
\newcommand{\Om}{\Omega_1^+}
\newcommand{\F}{\mathcal{F}}

\begin{document}

\begin{center}
{\Large Counterexample Packet}\\[0.3em]
{\large The Gaussian--triangle product is not extremal}\\[0.45em]
\textbf{Status:} candidate counterexample, likely valid\\
\textbf{Scope:} Question 3, second clause; the first clause remains open
\end{center}

\section*{Source and original question}

The source is Jaming--Matolcsi--R\'ev\'esz,
\emph{On the extremal rays of the cone of positive, positive definite
functions}, arXiv:0801.0941, published in the Journal of Fourier Analysis and
Applications \cite{JMR}.  Question 3 occurs in the Conclusion on PDF page 18.

\begin{center}
\includegraphics[width=\textwidth]{figures/open_problem_crop.png}
\end{center}

Writing
\[
 f=\mathbf 1_{[-1/2,1/2]}*\mathbf 1_{[-1/2,1/2]}
   =(1-|x|)_+,
 \qquad
 \gamma(x)=e^{-\pi x^2},
\]
the source asks whether $f*\gamma$ is extremal in $\Om$ and, if so, whether
\[
 G(x):=\gamma(x)(f*\gamma)(x)
\]
is extremal.  The result below gives a negative answer for $G$.

\section*{Definitions and Fourier convention}

We use
\[
 \widehat v(\xi)=\int_{\R}v(x)e^{-2\pi i x\xi}\,dx.
\]
The cone $\Om$ consists of the continuous even functions on $\R$ that are
both nonnegative and positive definite.  For the integrable functions used
below, nonnegativity of the Fourier transform implies positive definiteness
by Bochner's theorem and Fourier inversion \cite{Bochner,JMR}.

A nonzero $V\in\Om$ is an extremal ray generator if every decomposition
$V=V_1+V_2$ with $V_1,V_2\in\Om$ has $V_j=c_jV$ for nonnegative constants
$c_j$.

\section*{Proof intuition}

Completing the square moves half of the Gaussian from the pointwise product
into a convolution:
\[
 G=2q*\gamma_2,
 \qquad
 q(x)=(1-2|x|)_+e^{-2\pi x^2},
 \qquad
 \gamma_2(x)=e^{-2\pi x^2}.
\]
The compactly supported factor $q$ lies in $\Om$ and has a strictly positive
Fourier transform $Q$.  The jumps of $q'$ at $-1/2,0,1/2$ show that $Q$ is
bounded below by a positive multiple of $(1+\xi^2)^{-1}$.  Multiplying $q$
by $x^2$ produces a perturbation whose Fourier transform has the same upper
decay order.  Consequently, for sufficiently small $\varepsilon>0$, both
$q(1+\varepsilon x^2)$ and $q(1-\varepsilon x^2)$ remain nonnegative on both
the time and Fourier sides.  Convolving these two nonproportional functions
with $\gamma_2$ gives a nontrivial decomposition of $G$ in $\Om$.

\section*{Counterexample theorem}

\begin{theorem}
Let $f=(1-|x|)_+$ and $\gamma(x)=e^{-\pi x^2}$.  Then
\[
 G(x)=\gamma(x)(f*\gamma)(x)
\]
does not generate an extremal ray of $\Om$.

More precisely, with
\[
 q(x)=(1-2|x|)_+e^{-2\pi x^2},
 \qquad \gamma_2(x)=e^{-2\pi x^2},
\]
there is an $\varepsilon>0$ such that
\[
 q_\pm(x)=q(x)(1\pm\varepsilon x^2)
\]
both belong to $\Om$, and
\[
 G=(q_+*\gamma_2)+(q_-*\gamma_2)
\]
is a decomposition into two nonproportional elements of $\Om$.
\end{theorem}

\begin{proof}
\textbf{1. Gaussian refactorization.}
Since $f$ is supported on $[-1,1]$,
\begin{align*}
 G(x)
 &=e^{-\pi x^2}\int_{-1}^{1}f(t)e^{-\pi(x-t)^2}\,dt\\
 &=\int_{-1}^{1}f(t)e^{-2\pi(x-t/2)^2}e^{-\pi t^2/2}\,dt\\
 &=2\int_{-1/2}^{1/2}f(2s)e^{-2\pi s^2}e^{-2\pi(x-s)^2}\,ds\\
 &=2(q*\gamma_2)(x).
\end{align*}

\textbf{2. Strict positivity of $\widehat q$.}
Put $h(x)=(1-2|x|)_+=f(2x)$.  Both $h$ and $\gamma_2$ belong to $\Om$,
and $q=h\gamma_2$; hence $q\in\Om$.  Moreover,
\[
 \widehat h(\xi)
 =\frac12\left(\frac{\sin(\pi\xi/2)}{\pi\xi/2}\right)^2\geq0,
 \qquad
 \widehat{\gamma_2}(\xi)=\frac1{\sqrt2}e^{-\pi\xi^2/2}>0.
\]
Thus, if $Q=\widehat q$, then
\[
 Q=\widehat h*\widehat{\gamma_2}>0
 \qquad\hbox{on }\R,
\]
because the Gaussian factor is everywhere positive and $\widehat h$ is
nonnegative and nonzero away from only a discrete set.

\textbf{3. A global lower bound for $Q$.}
Extend $q$ by zero outside $[-1/2,1/2]$.  Its derivative has jumps
\[
 [q']_{-1/2}=2e^{-\pi/2},
 \qquad [q']_0=-4,
 \qquad [q']_{1/2}=2e^{-\pi/2}.
\]
Consequently, in the sense of distributions,
\[
 q''=r-4\delta_0
       +2e^{-\pi/2}(\delta_{-1/2}+\delta_{1/2}),
\]
where $r\in L^1(\R)$ is the ordinary second derivative on the smooth
pieces.  Taking Fourier transforms gives, for $\xi\ne0$,
\[
 -4\pi^2\xi^2Q(\xi)
 =\widehat r(\xi)-4+4e^{-\pi/2}\cos(\pi\xi).
\]
The Riemann--Lebesgue lemma yields
\[
 Q(\xi)
 =\frac{1-e^{-\pi/2}\cos(\pi\xi)}{\pi^2\xi^2}
  +o(\xi^{-2})
 \qquad (|\xi|\to\infty).
\]
Since $1-e^{-\pi/2}\cos(\pi\xi)\geq1-e^{-\pi/2}>0$, this gives a
positive lower bound $Q(\xi)\geq c_\infty/\xi^2$ outside a compact
interval.  On that compact interval, strict positivity and continuity of
$Q$ give another positive lower bound.  Therefore there is $c>0$ such that
\begin{equation}
 Q(\xi)\geq\frac{c}{1+\xi^2}
 \qquad(\xi\in\R).
 \tag{1}
\end{equation}

\textbf{4. A matched perturbation bound.}
Let $u(x)=x^2q(x)$ and $U=\widehat u$.  The function $u$ is continuous,
compactly supported, piecewise smooth, and its distributional second
derivative is a finite signed measure.  Hence
\[
 |U(\xi)|\leq\frac{C}{1+\xi^2}
 \qquad(\xi\in\R)
 \tag{2}
\]
for some $C>0$: for large $|\xi|$ this follows from
$-4\pi^2\xi^2U(\xi)=\widehat{u''}(\xi)$, and boundedness on a compact
interval follows from $u\in L^1$.

Choose
\[
 0<\varepsilon<\min\left\{4,\frac{c}{2C}\right\}.
\]
Because $|x|\leq1/2$ on the support of $q$,
$q_\pm=q(1\pm\varepsilon x^2)$ are nonnegative.  Equations (1)--(2) give
\[
 \widehat{q_\pm}(\xi)=Q(\xi)\pm\varepsilon U(\xi)
 \geq Q(\xi)-\varepsilon|U(\xi)|
 \geq\frac{c/2}{1+\xi^2}>0.
\]
The functions $q_\pm$ are continuous and integrable, and their Fourier
transforms are nonnegative (indeed integrable, by the same piecewise-smooth
decay argument).  Therefore $q_\pm\in\Om$.

\textbf{5. Nontrivial decomposition of $G$.}
Since $q_++q_-=2q$, Step 1 gives
\[
 (q_+*\gamma_2)+(q_-*\gamma_2)
 =2q*\gamma_2=G.
\]
Each summand belongs to $\Om$.  They are not proportional to $G$: if, for
example, $q_+*\gamma_2=\lambda G$, taking Fourier transforms and cancelling
the nowhere-zero transform of $\gamma_2$ would give
$q_+=2\lambda q$, contrary to $q_+/q=1+\varepsilon x^2$ on
$(-1/2,1/2)$.  Hence the displayed decomposition is nontrivial, and $G$ is
not extremal.
\end{proof}

\section*{Verification report}

\begin{center}
\begin{tabular}{>{\raggedright\arraybackslash}p{0.29\textwidth}>{\raggedright\arraybackslash}p{0.62\textwidth}}
\textbf{Step} & \textbf{Independent check}\\ \hline
Refactorization & Expanding
$2(x-t/2)^2+t^2/2$ gives $x^2+(x-t)^2$ exactly; the substitution
$t=2s$ supplies the factor $2$.\\
Fourier strict positivity & $q=h\gamma_2$ makes $Q$ a convolution of a
nonnegative sinc-square with a strictly positive Gaussian.\\
Jump asymptotic & Direct one-sided differentiation gives jumps
$2e^{-\pi/2},-4,2e^{-\pi/2}$; the resulting leading bracket is uniformly
positive because $e^{-\pi/2}<1$.\\
Perturbation & $u''$ is a finite measure, so $U=O(\xi^{-2})$, matching the
proved lower order of $Q$.  A single small $\varepsilon$ works globally.\\
Nonproportionality & Gaussian convolution is injective because its Fourier
transform never vanishes.\\
\end{tabular}
\end{center}

No computational step is used in the proof.  The auxiliary script
\texttt{code/numerical\_sanity\_check.py} checks the exact refactorization
and samples $Q\pm\widehat{x^2q}$ on $0\leq\xi\leq100$ for
$\varepsilon=1$.  Finite sampling is recorded only as a sanity check and is
not evidence for the global inequalities.

\section*{Scope and novelty check}

This theorem completely answers the second clause of Question 3 in the
negative.  It does not decide the first clause, namely whether $f*\gamma$ is
extremal.  It also shows that the function proposed in the source cannot
serve as a counterexample to Borisov's conjecture, because it is not
extremal; it does not decide that broader conjecture.

Before promotion, the run's registry, solution, attempt, and proof-gap
indexes were searched for arXiv:0801.0941, the source title, ``extremal
rays,'' and the Gaussian--triangle expressions.  No prior packet or attempt
was found.  External searches used the exact Question 3 phrases, the source
title and DOI, ``product of two extremals,'' ``positive positive-definite,''
``Borisov's conjecture,'' and arXiv-restricted variants.  They located the
source paper and general later work on Fourier positivity, but no paper
stating an answer to either exact clause.  Novelty confidence is therefore
moderate rather than absolute; citation databases and non-indexed literature
were not exhaustively searched.

\section*{Human review recommendation}

Send for mathematical review.  The highest-value checks are the
distributional jump computation in Step 3, the global lower-bound passage,
and the assertion that $u''$ is a finite measure in Step 4.  The construction
is otherwise elementary and has no unproved dependency.

\begin{thebibliography}{9}

\bibitem{JMR}
P. Jaming, M. Matolcsi, and S. Gy. R\'ev\'esz,
\emph{On the extremal rays of the cone of positive, positive definite
functions}, J. Fourier Anal. Appl. 15 (2009), 561--582;
arXiv:0801.0941,
\href{https://doi.org/10.1007/s00041-008-9057-6}{doi:10.1007/s00041-008-9057-6}.

\bibitem{Bochner}
S. Bochner,
\emph{Monotone Funktionen, Stieltjessche Integrale und harmonische Analyse},
Math. Ann. 108 (1933), 378--410.

\end{thebibliography}

\end{document}
